\SourceLeaf{206}Man kann die Darstellung also folgendermaßen fassen. Die Gleichung
\[
\left.\begin{array}{c} S \\ \Sigma \end{array}\right\} = M\pm h+i+u
\]
zeigt, dass gleiche Stellungen der Blase bei unverändertem Wert von $i$ gleiche Werte von $u$ anzeigen. Nun sind die Werte von $u$ in I und II entgegengesetzt. Hat [die Blase] also in II dieselbe Stellung wie in I, so ist $u'=u''=0$.

Im entgegengesetzten Fall schließt man, dass $u$ nicht $=0$ war und nun den entgegengesetzten Wert zum vorigen angenommen hat. Bringt man also durch stetiges Schrauben an 7 den Wert von $u$ auf den vorigen Wert zurück, so geht er dabei durch 0, und das halbe Zurückschrauben ergibt 0. Dann ist also die Linie 1\,2 oder 2\,1 oder 3\,4 horizontal, doch auch die Blase hat sich ihrer 2$^{\mathrm{ten}}$ Stellung wieder um die Hälfte genähert. Man ändert nun $i$ durch Schrauben an 3, bis $\left.\substack{S\\\Sigma}\right\} = M\pm h$ gilt.

\begin{center}
\textsc{Bemerkungen.}
\end{center}

Die erste der vorstehenden Notizen findet sich auf dem hinteren Schmutzblatt des Werkes \emph{Mathematische Abhandlungen} von Jacob Wilhelm Heinrich \textsc{Lehmann}, Zerbst 1829, und stammt vermutlich aus diesem Jahr. Die Notiz [2] steht auf einem einzelnen Blatt. Die Zeit ihrer Abfassung lässt sich nicht genau angeben, doch dürfte sie zwischen 1820 und 1830 liegen. Die Notiz [3] hat \textsc{Gauss} in einem Handbuch verzeichnet, höchstwahrscheinlich im März 1832, denn sie bezieht sich auf \textsc{Gauss}’ Brief an \textsc{Bolyai} vom 6. März 1832. Die Notiz [4] beruht auf einer mündlichen Mitteilung von \textsc{Gauss} an \textsc{Lübsen}, der im Jahr 1830 bei ihm Vorlesungen hörte (vgl. \textsc{Lübsen}, \emph{Ausf\"uhrliches Lehrbuch der Analysis}, Hamburg 1853, S. 171). Die Notiz [5] steht schließlich auf einem undatierten Zettel. Die Handschrift deutet jedoch auf eine spätere Periode in \textsc{Gauss}’ Leben, etwa auf die Jahre 1840--1850. Wahrscheinlich hat \textsc{Gauss} die Notiz bei der Vorbereitung einer Vorlesung über praktische Astronomie niedergeschrieben. Für ihre Aufnahme an dieser Stelle war allein ihr Inhalt maßgebend.

Dass \textsc{Gauss} sich schon sehr früh mit der Definition der Ebene beschäftigt hat, zeigt die bereits auf S. 162 dieses Bandes abgedruckte Stelle aus seinem Tagebuch vom 28. Juli 1797.

Die beiden Figuren in der Notiz [1] wurden dem Text hinzugefügt.

\hfill \textsc{Stäckel.}

